Shared symbols and units are listed in Table~\ref{tab:notation}. The following assumptions specify how we implement the benchmark \cite{problem}.

\textbf{Assumption 1 (Planar motion and rounded clearance).} Vehicles travel at constant speed $v=3$~m/s and occupy disks of radius $r_v=1.5$~m. The model omits acceleration and turning-radius limits. A vehicle center is admissible if it lies in the inset $[1.5,98.5]^2$ and at least $r_v$ from every building footprint \cite{lozano1983}. Direct clear legs use straight segments. Detours use an offset-corner visibility graph with continuous Euclidean clearance checked on every edge.

\textbf{Assumption 2 (Body tangency and optical occlusion).} Body tangency is permitted. A camera at $C=(x,y,h_c)$ with $h_c=2$~m can see a target $T$ only if the closed segment $\overline{CT}$ has empty intersection with all closed building prisms $B_k\times[0,H_k]$. This optical test excludes grazing contact with any face, edge, or vertex.

\textbf{Assumption 3 (Move, slew, observe, reset).} A vehicle begins its slew after it stops. The two axes rotate simultaneously from $(0^\circ,0^\circ)$ at up to $30^\circ$/s within the elevation interval $[-10^\circ,80^\circ]$. The axis that takes longer determines the slew time. Joint observation starts once both vehicles are ready at distinct non-collinear stations and lasts 5~s. Reuse or cancellation of a partially completed slew requires returning the gimbal to zero before moving.

\textbf{Assumption 4 (Measurement noise and drift).} Problem 2(1) assumes independent zero-mean Gaussian bearing errors with standard deviation $1^\circ$ and known altitude. Problem 3 assumes isotropic Gaussian position error with coordinate standard deviation $\sigma_{\mathrm{loc}}(s)=a+bs$ after traveling distance $s$. The independent gimbal offset and control errors have standard deviations $0.8^\circ$ and $0.6^\circ$. The combined horizontal error is wrapped onto $[-180^\circ,180^\circ)$.

\textbf{Assumption 5 (Event clocks and cached instructions).} Targets hover during the 60~s window. Disturbances occur at $t_d=5$~s. We take the specified absolute recovery clock $t_{\mathrm{rec}}\in(5,30]$~s as the baseline and test a 30~s outage sensitivity case at $t_{\mathrm{rec}}=35$~s. During an outage, vehicles follow cached nominal instructions. An ongoing-only case tests the effect of having no cached instructions for subsequent assignments.
